\chapter{Robustness, Security Limitations and Scalability Evaluation}
\label{chap:robustness-security-scalability}

The previous chapters presented the configuration model, the validation
pipeline, the generation of Kathar\'a artefacts, and the experimental use of
the framework. This chapter evaluates the framework from a different
perspective. Rather than considering only whether a valid laboratory can be
generated, it investigates how the translator behaves when its input becomes
large, structurally complex, ambiguous, or potentially unsafe.

From the perspective of the software-development lifecycle introduced in
Section~\ref{sec:software-development-lifecycle}, this chapter represents a
final verification stage. Earlier iterations primarily asked whether the
framework could generate and execute the intended laboratories. The tests below
instead challenge assumptions made during implementation and ask what happens
when the same mechanisms are exercised with large, malformed, ambiguous, or
potentially unsafe inputs. The resulting observations are therefore feedback on
the design choices described in Chapters~\ref{chap:requirements-architecture}--\ref{chap:laboratory-generation}, not only a list of independent limitations.

The analysis has two objectives. The first is to identify scalability limits
through reproducible stress tests. The second is to discuss implementation
choices that may become robustness or security weaknesses when configuration
files are malformed or originate from an untrusted source. These two aspects
are related. A configuration language is not only a convenient input format:
it is also an interface through which a user can request the creation of
thousands of nodes, routes, routing statements, files, and startup commands.

It is important to distinguish three categories throughout this chapter.
A \emph{security vulnerability} is a behaviour that can become exploitable
across an explicit trust boundary. A \emph{robustness weakness} is a condition
that can result in an incorrect or misleading laboratory without necessarily
creating a security compromise. A \emph{model limitation} is a deliberate or
implicit restriction in the set of networks that the translator can
represent. Not all limitations described below should therefore be interpreted
as exploitable vulnerabilities.

\section{Stress-Test Methodology}
\label{sec:stress-methodology}

The descriptive scalability benchmarks exercise the translator itself rather
than the runtime scalability of Docker or Kathar\'a. Each input is therefore processed
in two modes. The first uses the framework's validation-only option and
measures YAML loading, structural validation, and semantic validation without
creating a laboratory. The second performs complete generation and additionally
writes \texttt{lab.conf}, node startup files, and FRRouting files where
required. Starting the resulting containers is intentionally outside the scope
of these measurements because it would mix translator performance with Docker,
Kathar\'a, kernel, and host-resource limitations.

The benchmark data used in this chapter are the results recorded in
\path{final_benchmark.csv} and \path{final_summary.csv}. The CSV data
contain two measured repetitions for each available scenario and mode. No
measured execution reached a timeout. Every valid configuration returned exit
code 0 in both repetitions. The intentionally invalid configuration returned
exit code 1 in both repetitions, as expected. Because two repetitions cannot
support strong distributional or inferential claims, the measurements are used
only as descriptive observations of the tested executions.

The scenarios were selected so that different dimensions of input complexity
could be isolated rather than increasing only the number of nodes:

\begin{itemize}
    \item \texttt{flat\_10000.yml}: 10,000 nodes sharing a large flat
    topology, intended to stress node and address processing without creating
    thousands of distinct collision domains;
    \item \texttt{star\_4000.yml}: approximately 4,000 nodes arranged around
    a central router with a separate point-to-point collision domain for each
    peripheral connection, intended to stress subnet and collision-domain
    validation;
    \item \texttt{chain\_1000.yml}: 1,000 routers connected in a chain,
    stressing repeated routed-node and interface generation;
    \item \texttt{enterprise\_100r\_10h.yml}: 100 routers with ten hosts per
    routed LAN, providing a more realistic mixed topology;
    \item \texttt{dense\_500.yml}: 500 routers with multiple inter-router
    links, increasing interface density;
    \item \texttt{ospf\_1000.yml}: 1,000 OSPF routers, stressing FRRouting
    validation and configuration generation;
    \item \texttt{routes\_20000.yml}: a configuration dominated by 20,000
    static-route declarations;
    \item \texttt{commands\_100000.yml}: a configuration dominated by 100,000
    custom startup commands;
    \item \texttt{large\_yaml\_15mb.yml}: a 15.64 MiB YAML configuration,
    intended to stress deserialisation and in-memory representation;
    \item \texttt{invalid\_5000.yml}: 5,000 intentionally invalid nodes,
    intended to verify that large erroneous inputs are rejected without
    producing a laboratory.
\end{itemize}

For each measured run the benchmark records elapsed time, peak resident memory,
exit status, input size, number of generated files, and output size. Table
\ref{tab:stress-results} reports the aggregated values contained in the
provided summary CSV. Times are medians of the two recorded repetitions, while
memory is the maximum peak RSS observed in those repetitions.

\begin{table}[htbp]
\centering
\scriptsize
\resizebox{\textwidth}{!}{%
\begin{tabular}{lrrrrrrr}
\toprule
\textbf{Scenario} &
\textbf{Validate (s)} &
\textbf{Generate (s)} &
\textbf{Val. RSS (MiB)} &
\textbf{Gen. RSS (MiB)} &
\textbf{Input (MiB)} &
\textbf{Files} &
\textbf{Output (MiB)} \\
\midrule
Flat 10,000 nodes       & 4.065  & 8.840  & 97.04 & 109.16 & 1.127  & 10004 & 1.146 \\
Star 4,000 nodes        & 14.565 & 29.570 & 60.69 & 67.36  & 0.688  & 4004  & 0.727 \\
Chain 1,000 routers     & 1.320  & 2.835  & 27.03 & 28.66  & 0.165  & 4001  & 0.579 \\
Enterprise 100r/10h     & 0.535  & 1.095  & 25.53 & 26.96  & 0.141  & 1401  & 0.185 \\
Dense 500 routers       & 1.235  & 2.520  & 24.79 & 26.29  & 0.135  & 2001  & 0.347 \\
OSPF 1,000 routers      & 1.730  & 3.695  & 35.38 & 37.88  & 0.292  & 4001  & 0.675 \\
20,000 static routes    & 2.900  & 5.920  & 82.34 & 92.83  & 1.020  & 6     & 0.867 \\
100,000 commands        & 4.410  & 8.645  & 76.81 & 87.23  & 3.709  & 2     & 2.946 \\
Large YAML              & 7.610  & 15.240 & 74.47 & 124.36 & 15.640 & 2     & 15.030 \\
Invalid 5,000 nodes     & 2.025  & --     & 57.04 & --     & 0.595  & 0     & 0 \\
\bottomrule
\end{tabular}}
\caption{Descriptive stress-test results from the supplied benchmark CSV
files. Times are medians of two raw measurements and should not be interpreted
as stable population estimates.}
\label{tab:stress-results}
\end{table}

The results show that node count alone is not a sufficient predictor of
translation cost. The clearest example is the comparison between the flat and
star topologies. The flat topology contains 10,000 nodes and validates in
4.065 seconds. The star topology contains only about 4,000 nodes but requires
14.565 seconds for validation. It is therefore approximately 3.6 times slower
despite containing substantially fewer nodes. Complete generation exhibits the
same effect: 8.840 seconds for the 10,000-node flat case and 29.570 seconds for
the 4,000-node star case.

This difference is explained by the semantic checks performed on collision
domains. The validator compares the subnets associated with different
collision domains in order to reject overlapping Layer~3 address spaces.
When a topology contains a very large number of distinct point-to-point
domains, the number of pairwise comparisons can grow approximately as

\[
    \frac{D(D-1)}{2},
\]

where \(D\) is the number of collision domains being compared. In the
worst case this corresponds to quadratic growth, \(O(D^2)\). A flat topology,
by contrast, may contain many nodes while keeping \(D\) small. The stress
results therefore identify the number of distinct collision domains as a more
important scalability factor than raw node count for this part of semantic
validation.

The tests dominated by configuration volume reveal a different limit. The
20,000-route configuration validates in 2.900 seconds and reaches 82.34 MiB
peak RSS, while the 100,000-command configuration validates in 4.410 seconds
and reaches 76.81 MiB. The 15.64 MiB YAML file requires 7.610 seconds for
validation and reaches 74.47 MiB during validation. During complete generation
its peak RSS rises to 124.36 MiB, the highest value in the supplied summary.
These cases indicate that large nested YAML structures impose a cost even when
the logical topology is simple, because the complete document must be
deserialised into Python dictionaries and lists before semantic processing can
begin.

Another consistent result is the relationship between validation and complete
generation time. Across all valid scenarios, the ratio between the median
generation time and the median validation time ranges from approximately
1.96 to 2.17, with a mean close to 2.06. This behaviour is consistent with the
current implementation, in which normal execution validates the configuration
at the command-line boundary and the lower-level generation function loads and
validates it again before writing files. The duplicated validation gives
\texttt{generate\_lab()} an independent safety property, because the function
does not assume that its caller has already validated the input. However, the
stress tests show that this design also becomes measurable overhead for large
configurations.

Finally, \texttt{invalid\_5000.yml} provides a useful negative result. Both
recorded executions terminated in about 2.02--2.03 seconds with exit code 1,
without timeout and without generating output files. This demonstrates that
the validation pipeline can reject a large erroneous configuration without
continuing into the generation phase. From a robustness perspective this is a
positive property.

\section{Security and Robustness Weaknesses}
\label{sec:translator-weaknesses}

The stress tests identify performance boundaries, but other weaknesses arise
from the expressive features of the configuration language itself. Several of
these features are deliberate escape mechanisms. They are useful when YAML
files are trusted and authored by the laboratory designer, but they create a
clear trust boundary if configurations are downloaded, shared, or generated by
another user.

\paragraph{Arbitrary startup commands.}
The most important security-sensitive feature is the \texttt{commands} list.
Each command string is copied literally to the node startup file. This is
intentionally flexible: it allows the framework to launch applications,
configure kernel parameters, create files, or perform operations not covered
by the declarative model. It also means that a YAML file can carry executable
shell commands.

A harmless proof-of-concept configuration is:

\begin{lstlisting}[caption={Benign demonstration of command execution through a startup entry.},
label={lst:command-execution-example}]
nodes:
  pc_a:
    type: host
    interfaces:
      - network: lan_a
        ip: 10.0.1.10/24
    commands:
      - echo "COMMAND_EXECUTED" > /tmp/translator_test
\end{lstlisting}

The generated \texttt{pc\_a.startup} file contains the command literally. When
the laboratory starts, the command executes inside the node container. This
should not be described as an automatic host compromise: execution occurs in
the container context, and escalation beyond that boundary depends on the
container configuration and runtime privileges. Nevertheless, if YAML is
accepted from an untrusted source, the feature becomes a command-execution
boundary and should be treated accordingly.

A theoretical mitigation is to introduce two operating modes. A strict or
safe mode would reject \texttt{commands} entirely and permit only structured
operations represented by dedicated YAML fields. An explicit
\cliopt{allow-commands} option could retain the current behaviour for trusted
laboratories. Common operations currently expressed as shell strings could
progressively be converted into typed fields. Additional protection could
include command allowlists, clear warnings when a configuration contains
opaque commands, and reduced Linux capabilities for containers that do not
need privileged operations.

\paragraph{Opaque commands can bypass semantic intentions.}
The parser performs only limited consistency checks on command strings. For
example, it can recognise some explicit attempts to disable IPv4 forwarding on
a node declared as a router. General shell semantics, however, cannot be
reliably inferred from text. Equivalent operations can be written in many
forms. A configuration such as

\begin{lstlisting}[caption={Equivalent command that changes a forwarding runtime property.},
label={lst:forwarding-bypass-example}]
r1:
  type: router
  forwarding: true
  interfaces:
    - network: lan_a
      ip: 10.0.1.1/24
    - network: lan_b
      ip: 10.0.2.1/24
  commands:
    - sh -c 'echo 0 > /proc/sys/net/ipv4/ip_forward'
\end{lstlisting}

illustrates the general problem. The declarative field states that the node is
intended to forward traffic, while a later opaque command can change the
runtime state. In addition, \texttt{forwarding: true} currently expresses
intent and validation consistency but does not itself emit a kernel
\texttt{sysctl} operation.

The general solution is not to build an increasingly complex regular
expression for shell commands. A stronger design would make forwarding a
generated postcondition: when \texttt{forwarding: true} is selected, the
generator could explicitly configure the kernel state at a well-defined stage
and optionally verify it after startup. In strict mode, commands that directly
modify protected networking parameters could be disallowed.

\paragraph{Raw FRRouting configuration bypasses most semantic guarantees.}
The \texttt{frr.config} field provides another intentional escape mechanism.
Instead of generating protocol blocks from structured YAML fields, the
framework writes user-provided FRRouting text directly to
\texttt{frr.conf}. This is useful for advanced FRR features not represented by
the declarative model, but the translator cannot fully validate the semantics
of arbitrary FRR syntax.

For example:

\begin{lstlisting}[caption={Raw FRR content that cannot be fully validated by the YAML model.},
label={lst:raw-frr-example}]
frr:
  config: |
    router ospf
     this-command-does-not-exist
\end{lstlisting}

can be syntactically valid YAML while still containing an invalid FRR command.
Successful YAML translation therefore does not imply successful FRR startup.

A possible mitigation is to distinguish generated FRR configuration from raw
FRR configuration in the validation report and mark the latter as
``unverified''. A stricter implementation could invoke the FRR parser in a
disposable validation environment before accepting the laboratory. Another
option is to disable \texttt{frr.config} in strict mode and require the
structured protocol model whenever reproducibility and static validation are
more important than maximum expressiveness.

\paragraph{Static-route next hops are less deeply validated.}

Interfaces and default gateways receive strong semantic checks, including address parsing and directly connected gateway validation. Arbitrary static routes are less extensively analysed. Consider a router connecting two different networks:

\begin{lstlisting}[caption={Static route with a next hop outside the router's connected subnets.},
label={lst:static-route-weakness}]
r1:
  type: router
  interfaces:
    - network: lan_a
      ip: 10.0.1.1/24
    - network: lan_b
      ip: 10.0.2.1/24

  routes:
    - to: 172.16.0.0/16
      via: 192.168.200.1
\end{lstlisting}

The router is directly connected only to the \texttt{10.0.1.0/24} and \texttt{10.0.2.0/24} networks. However, the static route specifies \texttt{192.168.200.1} as its next hop, even though that address is not reachable through either configured interface.

The route can still be translated into:

\begin{lstlisting}[language=bash]
ip route add 172.16.0.0/16 via 192.168.200.1
\end{lstlisting}

Consequently, the YAML configuration may be accepted and translated even though the resulting route is operationally incorrect.

A possible improvement would be to apply to every \texttt{routes[].via} field a reachability check similar to the one already used for \texttt{default\_gateway}. The validator could verify that the next-hop address belongs to at least one subnet directly connected to the router. More advanced cases, such as explicitly configured \texttt{onlink} routes or recursive next-hop resolution, could later be supported through dedicated configuration options.

\paragraph{Unknown YAML keys can become silent misconfigurations.}
YAML itself permits arbitrary mapping keys. During rapid feature development,
retaining the loaded document as ordinary dictionaries made it easy to add new
properties without a separate schema-migration step. The same permissiveness
has a cost: the framework interprets the keys that belong to its model, while
an unknown node-level key is not automatically translated into Kathar\'a or
FRRouting behaviour. This creates a robustness risk when a user makes a
typographical error:

\begin{lstlisting}[caption={Typographical error that resembles a valid configuration field.},
label={lst:unknown-key-example}]
pc_a:
  type: host
  interfaces:
    - network: lan_a
      ip: 10.0.1.10/24
  defualt_gateway: 10.0.1.1
\end{lstlisting}

The user's intention is obvious, but \texttt{defualt\_gateway} is not
\texttt{default\_gateway}. A strict configuration schema should therefore
reject unknown fields rather than silently ignore them. This can be
implemented with explicit allowed-key sets or a schema-validation library.
Error messages could additionally suggest the closest valid field name.

\paragraph{FRR enable/disable semantics are not completely uniform.}
Router, switch, and firewall node types are considered FRR-capable by their
type, while a host becomes FRR-capable through its FRR declaration. As a
result, \texttt{frr.enabled: false} does not act as one universal
``disable FRR'' switch across all node types. This is primarily an API and
robustness weakness rather than a security vulnerability, because a
configuration may express a state whose operational interpretation is not
immediately intuitive.

A future version could define a single precedence rule: an explicit
\texttt{frr.enabled} value should override type-based defaults, while absence
of the field should preserve automatic capability detection. This would make
the meaning of the configuration independent of the path by which FRR
capability was inferred.

\paragraph{Container-image override introduces a trust decision.}
The optional \texttt{image} field overrides the default image selected by the
framework. This is necessary for custom laboratories, but a configuration from
an untrusted source can therefore request an image that has not been reviewed
by the laboratory author. The security consequences depend on Docker and
Kathar\'a runtime policy, so this should be described as an image-trust or
supply-chain risk rather than as a demonstrated container escape.

For controlled environments, image overrides could be limited to an allowlist
of repositories or immutable image digests. Strict mode could reject
unapproved images, while normal development mode would retain the current
flexibility.

\section{Translator Model Limitations}
\label{sec:translator-model-limitations}

Some behaviours that appear restrictive are consequences of the current
network model rather than vulnerabilities.

The validator assumes a single coherent Layer~3 address space for the complete
laboratory. IP addresses are treated as globally unique and overlapping
subnets on distinct collision domains are rejected. This simplifies target
resolution and catches many accidental addressing errors, but it prevents
legitimate multi-tenant or VRF-style laboratories in which isolated routing
domains intentionally reuse the same prefixes.

For example, the following conceptual topology is reasonable in a
multi-tenant system:

\begin{lstlisting}[caption={Address reuse that is meaningful only when routing domains are isolated.},
label={lst:overlapping-tenant-example}]
tenant_a_pc:
  interfaces:
    - network: tenant_a_lan
      ip: 192.168.1.10/24

tenant_b_pc:
  interfaces:
    - network: tenant_b_lan
      ip: 192.168.1.10/24
\end{lstlisting}

The current translator rejects this configuration because the same address is
used twice, even though the two logical tenants might never share a routing
domain. The corresponding mitigation would not be to weaken duplicate-address
validation globally. A better extension would introduce an explicit
\texttt{routing\_domain}, namespace, or VRF concept and scope address
uniqueness and subnet-overlap checks within that domain.

The switch and firewall node types have similar expressiveness limitations.
They reuse the common node and interface model, but the YAML language does not
currently provide a dedicated policy language for VLANs, trunks, spanning
tree, firewall rules, NAT policy, or complex filtering. Such behaviour must be
implemented through custom commands or imported node resources. This design is
sufficient for the current project scope but also increases reliance on the
opaque escape mechanisms discussed above.

\section{Possible Improvements}
\label{sec:possible-security-scalability-improvements}

The identified weaknesses can be addressed incrementally without redesigning
the complete framework.

First, the collision-domain overlap check should avoid an all-pairs comparison
when the number of domains becomes large. IPv4 and IPv6 networks can be
indexed separately and ordered by network address, or stored in a prefix or
interval structure. The objective is to reduce the common-case complexity from
quadratic pairwise comparison toward \(O(D \log D)\) indexing and lookup,
while preserving the current ability to report the two conflicting domains.

Second, the normal generation path could reuse the already loaded and
validated Python representation instead of parsing and validating the source a
second time. The lower-level generation function can retain its independent
safety property by exposing two interfaces: a public function that accepts a
path and performs complete validation, and an internal function that accepts a
validated configuration object. The command-line path would then use the
second function after its initial validation. The benchmark result that
generation takes approximately twice as long as validation across all
scenarios suggests that this change could have a measurable effect.

Third, explicit resource limits should be considered if configuration files
can originate from untrusted sources. Limits may include maximum YAML file
size, number of nodes, number of interfaces, number of collision domains,
number of static routes, number of custom commands, and maximum generated-file
count. These limits would convert extreme resource consumption into an early,
predictable validation error. The limits should be configurable so that
research users can deliberately raise them for large experiments.

Fourth, the configuration language could provide a strict mode. In this mode,
unknown fields would be errors, \texttt{commands} and raw
\texttt{frr.config} would be rejected or require explicit opt-in, image
overrides would be restricted, and structured route validation would be
stronger. The existing permissive mode could remain available for trusted
laboratory development.

Fifth, translation success could be separated from runtime correctness through
an optional post-generation verification stage. Such a stage could start the
laboratory, verify expected interface addresses and forwarding state, inspect
FRR daemon status, and run selected connectivity tests. This would not replace
static validation, but it would detect the class of errors that cannot be
proven by inspecting opaque shell or FRR commands.

\section{Discussion}
\label{sec:security-scalability-discussion}

The experiments show that the translator does not have a single meaningful
``maximum number of nodes''. Its effective scalability depends on which part
of the configuration grows. A 10,000-node flat topology validates in 4.065
seconds, while a 4,000-node star topology requires 14.565 seconds because it
contains many more independent collision domains that must be compared.
Large lists of routes and commands stress YAML deserialisation and memory
rather than topology analysis. Complete generation introduces additional
filesystem work and, in the current architecture, repeats loading and
validation.

The supplied results are encouraging in one respect: all valid stress
configurations completed without timeout or generation error, and the large
invalid configuration was rejected consistently without producing files.
The results therefore do not indicate a small hard capacity limit. Instead,
they reveal specific scaling dimensions that can be improved.

From a security perspective, the principal boundary is trust in the YAML
source. The framework already uses safe YAML loading, which prevents YAML tags
from directly constructing arbitrary Python objects. Safe deserialisation,
however, does not make the complete configuration harmless: the configuration
language deliberately contains fields that later become shell commands, raw
FRRouting text, and container-image choices. For a locally authored research
laboratory this is an acceptable and useful design decision. For a service
that accepted third-party configurations, the same features would require
strict sandboxing, explicit opt-in, resource limits, and stronger schema
validation.

The benchmark should also be interpreted with an important methodological
limitation. The supplied CSV records two measured repetitions per scenario.
The two runs are close to each other and are sufficient to identify large
differences such as the star-versus-flat behaviour, but they are not enough for
strong statistical claims about small performance differences. A future
evaluation could repeat each test five or more times on a documented reference
machine and report median, percentile, and variance. It would also be useful
to run a second, separate experiment that starts the generated laboratories in
Kathar\'a. That experiment would measure emulation scalability rather than
translator scalability and should not be mixed with the results presented
here.

Overall, the evaluation identifies a clear boundary between the parts of the
framework that are strongly validated and those deliberately left open for
advanced user control. The most relevant improvement areas are the
collision-domain overlap algorithm, duplicate loading and validation, strict
schema handling, static-route next-hop validation, and explicit controls for
opaque commands, raw FRR configuration, and custom images. These changes would
preserve the framework's main advantage---the ability to generate readable
standard Kathar\'a laboratories from a compact YAML source---while making its
behaviour more predictable under large or untrusted inputs.
